Проверка выполнена повторно для актуальной версии отчёта.
Скрипт \texttt{verify\_report.py} воспроизводит независимую квадратуру
коэффициентов и значения функции Грина. Для запуска требуются Python,
NumPy, SciPy и SymPy: \texttt{python verify\_report.py}.
\begin{center}
\begin{tabular}{lrr}
Тест & Параметры & Максимальная ошибка\\\hline
Стержень (12 мод) & $A=1.2$, $B=0.3$ & $1.39\cdot10^{-16}$\\
 & $b=0.7$, $T_e=0.4$, $l=a=1$ & \\
Шар (12 мод) & $q=k=c=R=1$ & $2.37\cdot10^{-14}$\\
Пластина (12 мод) & $Q=k=c=R=1$ & $1.12\cdot10^{-16}$
\end{tabular}
\end{center}
Числа округлены вверх; последние разряды могут зависеть от версий библиотек.
Для стержня сравниваются прямые интегралы
$2\int_0^1[x-U(x)]\sin(k_nx)dx$ и явные коэффициенты четвёртой темы.
Символьная подстановка $U$ в стационарное уравнение и обе границы
даёт тождественно нулевые невязки.
Для шара интегралы с весом $r^2$ сравниваются с
$D_n=-2/(z_n^2\cos z_n)$; ненулевые корни находятся методом Брента
по уравнению $\sin z-z\cos z=0$, без полюсов тангенса.
Для пластины проверяются коэффициенты $2(-1)^n/k_n^3$.
На сетке из 2001 точки при $N=100$, $l=a^2=1$, $b=0.5$, $\xi=0.4$
воспроизведены характеристики функции Грина:
\[
\begin{array}{c|ccc}
t&\max G_N&x_{\max}&\int_0^1G_Ndx\\\hline
0.01&2.8068780654&0.4000&0.9903580317\\
0.06&1.0555178537&0.4475&0.7295644344
\end{array}
\]
Интегралы здесь вычислены методом трапеций; координаты максимумов
относятся к выбранной сетке, а не к точной непрерывной оптимизации.
Скрипт проверяет согласованность формул, но не заменяет весь набор
компьютерных экспериментов и дополнительные вопросы преподавателя.
В среде повторной проверки XeLaTeX и LuaLaTeX отсутствовали:
фактическая сборка PDF и визуальная проверка вёрстки не выполнены.
Для сборки следует дважды выполнить \texttt{xelatex main.tex}
или \texttt{lualatex main.tex}.
